\subsection{Cartan行列の正定値性}

\begin{theorem}
  \label{thm:CartanMatrix.A_isPosDef}
  \lean{CartanMatrix.A_isPosDef}
  \leanok
  \uses{thm:sylvester_criterion, lem:SymmMatrix_leadingSubmatrix_comm, lem:CartanMatrix.A_leadingSubmatrix, thm:CartanMatrix.det_SymmMatrix_A}
  表~\ref{table:CartanMatrix.det_A}の各$A_n, B_n, \ldots, G_2$に対し，それを対称化した行列$S(A_n), S(B_n), \ldots, S(G_2)$はいずれも正定値である．
\end{theorem}

\begin{proof}
  例として$A_n$の場合を示す．
  定理~\ref{thm:sylvester_criterion} (Sylvesterの判定法)より，任意の主小行列式が正であることを言えば良い．
  補題~\ref{lem:SymmMatrix_leadingSubmatrix_comm}, \ref{lem:CartanMatrix.A_leadingSubmatrix}と定理~\ref{thm:CartanMatrix.det_SymmMatrix_A}より
  \begin{equation}
    \det (S(A_n))_k
    = \det S((A_n)_k)
    = \det S(A_k)
    = k + 1
    > 0
  \end{equation}
  となる．
  他も全く同様である．
\end{proof}

\begin{theorem}
  \label{thm:CCartanMatrix.A_tilda_isNotPosDef}
  \lean{CartanMatrix.A_tilda_isNotPosDef}
  \leanok
  \uses{thm:sylvester_criterion, thm:CartanMatrix.det_SymmMatrix_A_tilda}
  表~\ref{table:CartanMatrix.det_A_tilda}の各$\widetilde{A}_n, \widetilde{B}_n, \ldots, \widetilde{G}_2$に対し，それを対称化した行列$S(\widetilde{A}_n), S(\widetilde{B}_n), \ldots, S(\widetilde{G}_2)$はいずれも正定値でない．
\end{theorem}

\begin{proof}
  例として$\widetilde{A}_n$の場合を示す．
  定理~\ref{thm:sylvester_criterion} (Sylvesterの判定法)より，$k$次主小行列式が非正となる$k$が$1$つでもあればよい．
  $k = n + 1$とすると，定理~\ref{thm:CartanMatrix.det_SymmMatrix_A_tilda}より
  \begin{equation}
    \det (S(\widetilde{A}_n))_{n+1}
    = \det S(\widetilde{A}_n)
    = 0
    \le 0
  \end{equation}
  となる．
  他も全く同様である．
\end{proof}